\documentclass[12pt,a4paper]{article}

% ── Paquetes ──────────────────────────────────────────────────────────────────
\usepackage[utf8]{inputenc}
\usepackage[T1]{fontenc}
\usepackage[spanish]{babel}
\usepackage[margin=2.5cm]{geometry}
\usepackage{booktabs}
\usepackage[table]{xcolor}
\usepackage{fancyhdr}
\usepackage[many, breakable]{tcolorbox}
\usepackage{graphicx}
\usepackage{amssymb}

\definecolor{verdeosc}{RGB}{26,107,60}
\definecolor{verdemid}{RGB}{46,139,87}
\definecolor{dorado}{RGB}{212,172,13}
\definecolor{verdeclaro}{RGB}{232,244,232}

\setlength{\headheight}{30pt}
\pagestyle{fancy}
\fancyhf{}
\fancyhead[L]{\small\color{verdemid}SICVEC 2026}
\fancyhead[R]{\small\color{verdemid}Resumen -- Template}
\fancyfoot[C]{\small\thepage}
\renewcommand{\headrulewidth}{0.6pt}
\renewcommand{\footrulewidth}{0.4pt}

% ============================================================================
\begin{document}
% ============================================================================

\thispagestyle{empty}

\begin{center}
{\Large\bfseries\color{verdeosc} SIMPOSIO INTERNACIONAL DE CIENCIA VERDE Y ECONOMÍA CIRCULAR}\\
{\normalsize\color{verdemid} SICVEC 2026 -- 19 y 20 de octubre}\\[1cm]
{\Large\bfseries\color{verdeosc} TEMPLATE DE RESUMEN CIENTÍFICO}\\
\end{center}

\vspace{1cm}

\noindent
\begin{tcolorbox}[enhanced, colback=verdeclaro, colframe=verdeosc, boxrule=1.5pt, left=10pt, right=10pt, top=10pt, bottom=10pt]

\noindent
\textbf{\Large Título del Trabajo}

[El título debe ser conciso (máximo 15 palabras), descriptivo y atractivo.\\
Ejemplo: ``Evaluación de la Economía Circular en Agroindustrias Colombianas: Caso de Estudio en Producción de Biodiesel'']

\vspace{0.5cm}

\noindent
\textbf{Autores:} Nombre Completo$^{1*}$, Segundo Autor$^{2}$, Tercer Autor$^{3}$

\noindent
\textbf{Afiliaciones:}\\
$^{1}$ Departamento de Biología y Química, Universidad de Sucre, Sincelejo, Colombia\\
$^{2}$ Instituto de Investigación, Institución X, Ciudad, País\\
$^{3}$ Empresa de Tecnología Verde, Ciudad, País\\
$^{*}$Autor correspondiente: email@example.com, +XX XXXXXXXXX

\vspace{0.5cm}

\noindent
\textbf{Eje Temático:} Seleccionar uno:
\begin{itemize}
  \item $\square$ Economía Circular y Procesos Lineales
  \item $\square$ Biotecnología Sostenible
  \item $\square$ Tecnologías Verdes y Energías Renovables
  \item $\square$ Sostenibilidad en Cadenas de Suministro
  \item $\square$ Política Pública y Gobernanza Ambiental
\end{itemize}

\noindent
\textbf{Modalidad:} $\square$ Ponencia Oral \quad $\square$ Poster \quad $\square$ Taller

\end{tcolorbox}

\newpage

% ============================================================================
\section*{CONTENIDO DEL RESUMEN}
% ============================================================================

\vspace{1cm}

\noindent
\textbf{[SECCIÓN 1: INTRODUCCIÓN/JUSTIFICACIÓN]}\\
(Explicar por qué es importante este tema. Contexto global y local. Máx. 60 palabras)

\vspace{0.3cm}

La economía circular es fundamental para reducir la presión sobre recursos naturales. En Colombia, el sector agroindustrial genera aproximadamente 45 millones de toneladas de residuos anuales, de los cuales menos del 5\% se aprovecha. Este trabajo busca...

\vspace{0.8cm}

\noindent
\textbf{[SECCIÓN 2: OBJETIVOS]}\\
(¿Qué se pretende lograr? Objetivo general y 1-2 específicos. Máx. 50 palabras)

\vspace{0.3cm}

\textbf{General:} Evaluar la viabilidad de implementar modelos de economía circular en agroindustrias colombianas.

\textbf{Específicos:}
\begin{itemize}
  \item Identificar corrientes de residuos aprovechables en 5 plantas de procesamiento
  \item Proponer estrategias de simbiosis industrial
\end{itemize}

\vspace{0.8cm}

\noindent
\textbf{[SECCIÓN 3: METODOLOGÍA]}\\
(¿Cómo se desarrolló la investigación? Métodos, materiales, población estudiada. Máx. 60 palabras)

\vspace{0.3cm}

Se realizó un estudio de caso en 5 agroindustrias del Caribe colombiano. Se aplicaron auditorías ambientales cuantitativas de residuos (balance de masa), análisis de ciclo de vida (LCA) usando SimaPro 9.1, y entrevistas semiestructuradas con 15 gerentes de producción. Se empleó modelado de simbiosis industrial con software Umberto...

\vspace{0.8cm}

\noindent
\textbf{[SECCIÓN 4: RESULTADOS/HALLAZGOS]}\\
(¿Qué se encontró? Datos, cifras, descubrimientos principales. Máx. 80 palabras)

\vspace{0.3cm}

El 68\% de los residuos generados tienen potencial de revalorización. Se identificaron 12 oportunidades de simbiosis industrial, especialmente en aprovechamiento de biomasa para bioenergía (1.2 MW potenciales) y biofertilizantes. Una planta piloto de biodigestión anaeróbica generó un retorno de inversión en 2.3 años. Las barreras principales fueron técnicas (falta de tecnología local) y políticas (ausencia de incentivos verdes)...

\vspace{0.8cm}

\noindent
\textbf{[SECCIÓN 5: CONCLUSIONES/IMPLICACIONES]}\\
(¿Cuál es el impacto? Aportes a sostenibilidad y posibles aplicaciones. Máx. 60 palabras)

\vspace{0.3cm}

La economía circular es viable en agroindustrias colombianas con factibilidad económica comprobada. Se requiere: (1) transferencia de tecnología, (2) políticas públicas de incentivo, y (3) capacitación de operarios. Los resultados pueden replicarse en contextos similares de América Latina, contribuyendo a circularidad y descarbonización del sector...

\newpage

% ============================================================================
\section*{INFORMACIÓN COMPLEMENTARIA}
% ============================================================================

\vspace{0.8cm}

\noindent
\textbf{Palabras clave:} (3-5 palabras, separadas por comas, que resuman el tema)

\vspace{0.3cm}

Economía circular, simbiosis industrial, agroindustria, residuos, revalorización

\vspace{0.8cm}

\noindent
\textbf{Referencias Citadas:} (Máximo 5, en formato APA)

\vspace{0.3cm}

\begin{enumerate}
  \item Ellen MacArthur Foundation. (2023). Towards the Circular Economy: Economic and business rationale for an accelerated transition. www.ellenmacarthurfoundation.org

  \item Kirchherr, J., Reike, D., \& Hekkert, M. (2017). Conceptualizing the circular economy: An analysis of 114 definitions. Resources, Conservation and Recycling, 127, 221-232.

  \item Ministerio de Ambiente de Colombia. (2022). Estrategia Nacional de Economía Circular 2030.

  \item Lent, P. J. (2021). Circular Economy and Agricultural Residues in Latin America. Journal of Cleaner Production, 45(3), 234-245.

  \item Università di Verona (2024). Industrial Symbiosis Database - Latin America Region.
\end{enumerate}

\vspace{1cm}

% ============================================================================
\section*{CONSIDERACIONES FINALES PARA LOS AUTORES}
% ============================================================================

\begin{tcolorbox}[enhanced, colback=yellow!10, colframe=orange, boxrule=1pt, left=10pt, right=10pt, top=10pt, bottom=10pt]

\textbf{Verificar antes de enviar:}

\begin{itemize}
  \item[$\square$] Resumen total entre 250-300 palabras (excluye referencias)
  \item[$\square$] Estructura clara: Introducción → Objetivos → Metodología → Resultados → Conclusiones
  \item[$\square$] Título atractivo y descriptivo (máx. 15 palabras)
  \item[$\square$] Todos los autores con afiliaciones completas
  \item[$\square$] Email de autor correspondiente incluido
  \item[$\square$] Eje temático seleccionado correctamente
  \item[$\square$] Palabras clave relevantes (3-5)
  \item[$\square$] Sin errores de ortografía y gramática
  \item[$\square$] Máximo 5 referencias bibliográficas
  \item[$\square$] Formato: Times New Roman, 11pt, espaciado 1.5
  \item[$\square$] Texto del resumen sin nombres de autores ni afiliaciones (revisión ciega)
  \item[$\square$] Resumen copiado en el formulario en línea (no se sube archivo)
\end{itemize}

\end{tcolorbox}

\vspace{1.5cm}

\begin{center}
{\small\color{verdemid} Para enviar tu resumen, ingresa a la plataforma de inscripción:\\
\textbf{www.sicvec2026.unisucre.edu.co}}

{\small\color{verdeosc} Plazo: 4 de octubre de 2026}
\end{center}

\end{document}
